\scanpage{9}
seules droites invariantes \inserted{\unsure{pour|par} $\mathbb D_\infty$} de $\operatorname{Sym}^2(\check E)$, \deleted{les} celles qui sont \unsure{formées|fournies} de formes quadratiques \uline{invariantes}.

Si on permet \deleted{que} ${\unsure{\zeta|\xi}}^3=1$ sur certaines fibres, \margin{i.e. $\alpha=-1$ sur certaines fibres} \deleted{alors on} mais pas $\zeta^4=1$, alors on trouve que les droites inv. dans $\operatorname{Sym}^2(\check E)$ pour $u$ sont contenues \inserted{loc.} dans $V=\check D_1^{\otimes2}\unsure{\oplus}\check D_2$ ou $\unsure{W|V}=\check D_2^{\otimes2}\oplus\check D_1$ ou $\check D_0^{\otimes2}\oplus\check D_1\otimes\check D_2$, or $\sigma$ échange $V$ et $W$, donc il n'y a pas de droite inv. par $\Phi$ ni dans $V$, ni dans $W$, et on \deleted{\unsure{est dans les}}\inserted{obtient la} $=$ conclusion.

Mais si on permet $\zeta^4=1$ sur certaines fibres, que se passe-t-il ? Supposons donc carrément \deleted{$\zeta^4$}
\[
\zeta^4=1\qquad\text{i.e.}\qquad\zeta^2=-1,\ \text{ou encore }\alpha=0
\tag*{(4)}
\]
(cas du carré) -- car (4) s'écrit
\[
(\zeta^2-1)(\zeta^2+1)=0
\]
et comme par hyp. $\zeta^2-1$ est inv, celà donne $\zeta^2+1=0$ i.e. $\zeta+\zeta^{-1}=0$ i.e. $\alpha=0\ldots$

Alors les valeurs propres distinctes parmi (3) sont $\zeta^2=\zeta^{-2}=-1,\ 1,\ \zeta,\zeta^{-1}$, correspondant aux sous-espaces propres $\check D_1^{\otimes2}\oplus\check D_2^{\otimes2},\check D_0^{\otimes2}\oplus(D_1\otimes D_2),\check D_1,\check D_2$. Les droites propres
% archive sheet number: 8
